% !TEX root = main.tex
% Chương 3: Xây dựng đề cương nghiên cứu chi tiết (45% tổng điểm, 9-14 trang)
% Lưu ý kỹ thuật: mọi bảng/hình/công thức PHẢI được tham chiếu trong văn bản
% bằng \ref{...} -- không để bảng/hình/công thức "mồ côi".

\chapter{Xây dựng đề cương nghiên cứu chi tiết}
\label{ch:chuong3}

\begin{spacing}{1.2}

\section*{Giới thiệu chương}

\begin{huongdan}
Trọng số: 45\% tổng điểm, chương trọng tâm nhất của tiểu luận. Độ dài khuyến nghị: 10--15 trang. Thiết kế một đề cương nghiên cứu thực sự khả thi, đủ chi tiết để người đọc hình dung rõ cách thực hiện, đánh giá và triển khai đề tài. Mọi khẳng định phải được hỗ trợ bởi bằng chứng khoa học (trích dẫn từ Chương 2) hoặc số liệu thực tế từ doanh nghiệp, vì đây là nơi Chương 1, 2 và 3 phải liên kết chặt chẽ với nhau. Đoạn mở đầu này nên nêu ngắn gọn chương gồm những mục nào (Mục~\ref{sec:tieu_de_muc_tieu} đến Mục~\ref{sec:ket_qua_du_kien}) và mạch triển khai giữa chúng. Sau chương này cần đạt được: xây dựng tiêu đề chuẩn và các mục tiêu nghiên cứu đo lường được; kiểm tra tính khả thi qua khung FINER với bằng chứng cụ thể; thiết kế phương pháp nghiên cứu đủ chi tiết để triển khai thực tế; lập kế hoạch thực hiện, quản lý rủi ro và dự kiến kết quả đầu ra.
\end{huongdan}

Nội dung của phần này.

\section{Tiêu đề và mục tiêu nghiên cứu}
\label{sec:tieu_de_muc_tieu}

\subsection{Xây dựng tiêu đề nghiên cứu}
\label{subsec:tieu_de}

\begin{huongdan}
Một tiêu đề tốt phải trả lời được ba câu hỏi: phương pháp/công nghệ gì, để làm gì, cho ai/ở đâu. Trình bày tiêu đề chính thức; phân tích từng thành phần của tiêu đề; giải thích tiêu đề phản ánh đúng Research Gap đã xác định ở Chương 2 như thế nào; nêu ít nhất một tiêu đề thay thế đã cân nhắc và lý do không chọn.
\end{huongdan}

Nội dung của phần này.

\subsection{Mục tiêu nghiên cứu}
\label{subsec:muc_tieu}

\begin{huongdan}
Mục tiêu tổng quát (GO) trình bày trong một đến hai câu, bắt đầu bằng động từ hành động. Mục tiêu cụ thể (SO1, SO2...) gồm ba đến năm mục tiêu nhỏ, mỗi mục cụ thể, đo lường được và đóng góp trực tiếp cho mục tiêu tổng quát.
\end{huongdan}

Nội dung của phần này. Hệ thống mục tiêu nghiên cứu được trình bày tại Bảng~\ref{tab:muc_tieu}.

\begin{table}[htpb]
\centering
\caption{Mục tiêu nghiên cứu tổng quát và cụ thể}
\label{tab:muc_tieu}
\renewcommand{\arraystretch}{1.3}
{\fontsize{12}{14}\selectfont
\begin{tabularx}{\textwidth}{|>{\raggedright\arraybackslash}p{2.6cm}|X|>{\raggedright\arraybackslash}p{4cm}|}
\hline
\textbf{Mục tiêu} & \textbf{Nội dung} & \textbf{Chỉ số đo lường} \\
\hline
GO (tổng quát) & [Học viên điền] & N/A -- bao quát \\
\hline
SO1 & [Học viên điền] & [Chỉ số cụ thể] \\
\hline
SO2 & [Học viên điền] & [Chỉ số cụ thể] \\
\hline
SO3 & [Học viên điền] & [Chỉ số cụ thể] \\
\hline
\end{tabularx}
}
\end{table}

\section{Câu hỏi nghiên cứu và kiểm tra FINER}
\label{sec:rq_finer}

\subsection{Câu hỏi nghiên cứu (Research Questions)}
\label{subsec:rq}

\begin{huongdan}
Nêu một câu hỏi nghiên cứu chính (MRQ) bao quát toàn bộ đề tài, và ba đến năm câu hỏi phụ (RQ1, RQ2...) cụ thể, mỗi câu trả lời được bằng dữ liệu thu thập được. Lập bảng liên kết mỗi câu hỏi với phương pháp và kết quả tương ứng.
\end{huongdan}

Nội dung của phần này. Hệ thống câu hỏi nghiên cứu được tổng hợp tại Bảng~\ref{tab:rq}.

\begin{table}[htpb]
\centering
\caption{Câu hỏi nghiên cứu và phương pháp trả lời}
\label{tab:rq}
\renewcommand{\arraystretch}{1.3}
{\fontsize{12}{14}\selectfont
\begin{tabularx}{\textwidth}{|>{\raggedright\arraybackslash}p{1.8cm}|X|>{\raggedright\arraybackslash}p{3.6cm}|>{\raggedright\arraybackslash}p{3.2cm}|}
\hline
\textbf{Câu hỏi} & \textbf{Nội dung} & \textbf{Phương pháp trả lời} & \textbf{Kết quả kỳ vọng} \\
\hline
MRQ & [Học viên điền] & Toàn bộ quy trình & [Điền] \\
\hline
RQ1 & [Học viên điền] & [Điền] & [Điền] \\
\hline
RQ2 & [Học viên điền] & [Điền] & [Điền] \\
\hline
RQ3 & [Học viên điền] & [Điền] & [Điền] \\
\hline
\end{tabularx}
}
\end{table}

\subsection{Kiểm tra tiêu chí FINER}
\label{subsec:finer}

\begin{huongdan}
Trả lời từng tiêu chí (Feasible, Interesting, Novel, Ethical, Relevant) bằng bằng chứng cụ thể, không chỉ ghi "Có" hay "Đạt". Tiêu chí Novel phải liên kết với Research Gap ở mục~\ref{subsec:research_gap}; tiêu chí Relevant phải liên kết với Problem Statement ở mục~\ref{subsec:problem_statement}.
\end{huongdan}

Nội dung của phần này. Kết quả kiểm tra tính khả thi của đề tài theo khung \tvt{FINER}{Feasible, Interesting, Novel, Ethical, Relevant} được trình bày tại Bảng~\ref{tab:finer}.

\begin{table}[htpb]
\centering
\caption{Kiểm tra đề tài theo khung FINER}
\label{tab:finer}
\renewcommand{\arraystretch}{1.3}
{\fontsize{12}{14}\selectfont
\begin{tabularx}{\textwidth}{|>{\raggedright\arraybackslash}p{3.4cm}|>{\raggedright\arraybackslash}p{5.2cm}|X|}
\hline
\textbf{Tiêu chí} & \textbf{Câu hỏi đánh giá} & \textbf{Bằng chứng cụ thể} \\
\hline
F -- Feasible (Khả thi) & Có đủ dữ liệu, kỹ năng, ngân sách, thời gian? & [Học viên điền] \\
\hline
I -- Interesting (Thú vị) & Kết quả có ý nghĩa với khoa học và xã hội? & [Học viên điền] \\
\hline
N -- Novel (Mới mẻ) & Có gì khác so với nghiên cứu đã có? & [Liên kết Research Gap, mục~\ref{subsec:research_gap}] \\
\hline
E -- Ethical (Đạo đức) & Dữ liệu hợp pháp? Quyền riêng tư? & [Học viên điền] \\
\hline
R -- Relevant (Liên quan) & Giải quyết vấn đề thực tiễn cấp thiết? & [Liên kết Problem Statement, mục~\ref{subsec:problem_statement}] \\
\hline
\end{tabularx}
}
\end{table}

\section{Phương pháp nghiên cứu dự kiến}
\label{sec:phuong_phap_nc}

\subsection{Loại hình nghiên cứu và thiết kế tổng thể}
\label{subsec:thiet_ke}

\begin{huongdan}
Xác định loại hình nghiên cứu (thực nghiệm, khảo sát, case study, hay kết hợp) kèm lý do chọn; mô tả thiết kế nghiên cứu tổng thể bằng sơ đồ luồng; xác định biến độc lập, biến phụ thuộc, biến kiểm soát; nêu phương pháp đánh giá, kiểm chứng và cơ sở so sánh (baseline).
\end{huongdan}

Nội dung của phần này. Thiết kế nghiên cứu tổng thể được minh họa tại Hình~\ref{fig:thiet_ke_nc}.

\begin{figure}[htpb]
\centering
\begin{tikzpicture}[
  block/.style={rectangle, draw=green!60!black, fill=green!5, rounded corners,
                minimum height=1.1cm, text width=2.3cm, align=center, font=\small},
  arrow/.style={-{Stealth[scale=1.1]}, thick, draw=green!50!black}
]
\node[block] (b1) {Thu thập dữ liệu};
\node[block, right=0.55cm of b1] (b2) {Tiền xử lý};
\node[block, right=0.55cm of b2] (b3) {Xây dựng giải pháp};
\node[block, right=0.55cm of b3] (b4) {Đánh giá};
\node[block, right=0.55cm of b4] (b5) {Triển khai};
\draw[arrow] (b1) -- (b2);
\draw[arrow] (b2) -- (b3);
\draw[arrow] (b3) -- (b4);
\draw[arrow] (b4) -- (b5);
\end{tikzpicture}
\caption{Sơ đồ thiết kế nghiên cứu tổng thể (học viên điều chỉnh theo đề tài)}
\label{fig:thiet_ke_nc}
\end{figure}

\subsection{Quy trình thu thập dữ liệu}
\label{subsec:thu_thap_du_lieu}

\begin{huongdan}
Nêu nguồn dữ liệu (hệ thống, loại, định dạng, thời gian); phương pháp thu thập (truy vấn SQL, gọi API, xuất tệp, khảo sát...); quy mô dữ liệu dự kiến; quy trình làm sạch, chuẩn hóa, tạo đặc trưng; và kế hoạch xử lý khi gặp vấn đề dữ liệu.
\end{huongdan}

Nội dung của phần này.

\subsection{Phương pháp phân tích và công nghệ sử dụng}
\label{subsec:phan_tich_cong_nghe}

\begin{huongdan}
Mô tả thuật toán/mô hình sẽ áp dụng (nguyên lý hoạt động, tham số chính, lý do phù hợp); liệt kê đầy đủ công nghệ sử dụng (ngôn ngữ lập trình, thư viện, framework, cơ sở dữ liệu, nền tảng đám mây, công cụ trực quan hóa); nêu quy trình huấn luyện/thực hiện và kiến trúc hệ thống.
\end{huongdan}

Để minh họa cách trình bày một độ đo định lượng có công thức và được tham chiếu xuyên suốt bài viết, độ chính xác (Accuracy) của mô hình dự báo được định nghĩa tại Công thức~\eqref{eq:accuracy_vd}.

\begin{equation}
\text{Accuracy} = \frac{\text{Số dự báo đúng}}{\text{Tổng số dự báo}} \times 100\%
\label{eq:accuracy_vd}
\end{equation}

Sai số phần trăm tuyệt đối trung bình (\tvt{MAPE}{Mean Absolute Percentage Error}), chỉ số được dùng làm mục tiêu chính trong Bảng~\ref{tab:kq_du_kien}, được tính theo Công thức~\eqref{eq:mape_vd}, trong đó $A_t$ là giá trị thực tế, $F_t$ là giá trị dự báo tại thời điểm $t$, và $n$ là tổng số quan sát.

\begin{equation}
\text{MAPE} = \frac{100\%}{n} \sum_{t=1}^{n} \left| \frac{A_t - F_t}{A_t} \right|
\label{eq:mape_vd}
\end{equation}

Học viên thay hai công thức trên bằng độ đo thực sự phù hợp với đề tài của mình (ví dụ: Precision, Recall, F1-score, RMSE...), miễn là mỗi công thức có nhãn (\texttt{\textbackslash label}) và được tham chiếu ít nhất một lần trong văn bản bằng \texttt{\textbackslash eqref} hoặc \texttt{Công thức~\textbackslash ref\{...\}}, đúng như Công thức~\eqref{eq:accuracy_vd} và Công thức~\eqref{eq:mape_vd} đã minh họa.

Nội dung của phần này. Mô hình LSTM được lựa chọn vì đã cho thấy hiệu quả trong các bài toán dự báo chuỗi thời gian tương tự \cite{vidu_tapchi}, và được cài đặt bằng TensorFlow/Keras theo tài liệu chính thức của framework \cite{vidu_website}.

\section{Kế hoạch thực hiện}
\label{sec:ke_hoach}

\subsection{Bảng kế hoạch theo giai đoạn (Timeline)}
\label{subsec:timeline}

\begin{huongdan}
Chia rõ ít nhất năm giai đoạn với thời gian cụ thể; mỗi giai đoạn có tên, thời gian, hoạt động, kết quả đầu ra đo lường được và mốc quan trọng (Milestone). Khung thời gian tổng thể phải thực tế.
\end{huongdan}

Nội dung của phần này. Kế hoạch thực hiện theo giai đoạn được trình bày tại Bảng~\ref{tab:timeline}.

\begin{table}[htpb]
\centering
\caption{Kế hoạch thực hiện theo giai đoạn}
\label{tab:timeline}
\renewcommand{\arraystretch}{1.3}
{\fontsize{12}{14}\selectfont
\begin{tabularx}{\textwidth}{|>{\raggedright\arraybackslash}p{4.2cm}|>{\centering\arraybackslash}p{2.2cm}|X|>{\raggedright\arraybackslash}p{3.4cm}|}
\hline
\textbf{Giai đoạn} & \textbf{Thời gian} & \textbf{Hoạt động chi tiết} & \textbf{Kết quả đầu ra} \\
\hline
GĐ 1: Chuẩn bị và khảo sát & [Điền] & [Học viên điền] & [Điền] \\
\hline
GĐ 2: Thu thập và xử lý dữ liệu & [Điền] & [Học viên điền] & [Điền] \\
\hline
GĐ 3: Xây dựng giải pháp & [Điền] & [Học viên điền] & [Điền] \\
\hline
GĐ 4: Triển khai và kiểm thử & [Điền] & [Học viên điền] & [Điền] \\
\hline
GĐ 5: Viết báo cáo và thuyết trình & [Điền] & [Học viên điền] & [Điền] \\
\hline
\end{tabularx}
}
\end{table}

\subsection{Quản lý rủi ro (Risk Management)}
\label{subsec:rui_ro}

\begin{huongdan}
Xác định ít nhất năm rủi ro tiềm ẩn; đánh giá xác suất (thấp/trung bình/cao) và mức độ ảnh hưởng; mỗi rủi ro có phương án ứng phó cụ thể, không viết chung chung kiểu "sẽ xử lý sau".
\end{huongdan}

Nội dung của phần này. Các rủi ro tiềm ẩn và phương án ứng phó được tổng hợp tại Bảng~\ref{tab:rui_ro}.

\begin{table}[htpb]
\centering
\caption{Nhận diện rủi ro và phương án ứng phó}
\label{tab:rui_ro}
\renewcommand{\arraystretch}{1.3}
{\fontsize{12}{14}\selectfont
\begin{tabularx}{\textwidth}{|X|>{\centering\arraybackslash}p{2.7cm}|>{\centering\arraybackslash}p{2.7cm}|X|}
\hline
\textbf{Rủi ro} & \textbf{Xác suất} & \textbf{Ảnh hưởng} & \textbf{Phương án ứng phó cụ thể} \\
\hline
[Học viên điền] & [Thấp/TB/Cao] & [Thấp/TB/Cao] & [Học viên điền] \\
\hline
[Học viên điền] & [Thấp/TB/Cao] & [Thấp/TB/Cao] & [Học viên điền] \\
\hline
[Học viên điền] & [Thấp/TB/Cao] & [Thấp/TB/Cao] & [Học viên điền] \\
\hline
[Học viên điền] & [Thấp/TB/Cao] & [Thấp/TB/Cao] & [Học viên điền] \\
\hline
[Học viên điền] & [Thấp/TB/Cao] & [Thấp/TB/Cao] & [Học viên điền] \\
\hline
\end{tabularx}
}
\end{table}

\section{Kết quả dự kiến và ý nghĩa nghiên cứu}
\label{sec:ket_qua_du_kien}

\subsection{Kết quả kỹ thuật dự kiến}
\label{subsec:kq_ky_thuat}

\begin{huongdan}
Liệt kê các chỉ số kỹ thuật mục tiêu, cụ thể và đo lường được, so sánh với baseline hiện tại; giải thích cơ sở khoa học cho mỗi mục tiêu bằng cách liên kết với bài báo ở Chương 2; mô tả các sản phẩm đầu ra cụ thể (mô hình, API, dashboard, tài liệu kỹ thuật).
\end{huongdan}

Nội dung của phần này. Các chỉ số kỹ thuật mục tiêu được trình bày tại Bảng~\ref{tab:kq_du_kien}, trong đó chỉ số MAPE được tính theo Công thức~\eqref{eq:mape_vd} đã trình bày tại mục~\ref{subsec:phan_tich_cong_nghe}; các ngưỡng mục tiêu được đối chiếu với báo cáo kỹ thuật tham khảo \cite{vidu_baocao} để đảm bảo tính khả thi.

\begin{table}[htpb]
\centering
\caption{Kết quả kỹ thuật dự kiến}
\label{tab:kq_du_kien}
\renewcommand{\arraystretch}{1.3}
{\fontsize{12}{14}\selectfont
\begin{tabularx}{\textwidth}{|X|>{\centering\arraybackslash}p{3cm}|>{\centering\arraybackslash}p{2.6cm}|>{\raggedright\arraybackslash}p{4cm}|}
\hline
\textbf{Chỉ số} & \textbf{Hiện trạng (baseline)} & \textbf{Mục tiêu (target)} & \textbf{Cơ sở khoa học} \\
\hline
[Học viên điền] & [Điền] & [Điền] & [Trích dẫn từ Chương 2] \\
\hline
[Học viên điền] & [Điền] & [Điền] & [Trích dẫn từ Chương 2] \\
\hline
[Học viên điền] & [Điền] & [Điền] & [Trích dẫn từ Chương 2] \\
\hline
\end{tabularx}
}
\end{table}

\subsection{Tác động kinh doanh (Business Impact)}
\label{subsec:business_impact}

\begin{huongdan}
Lượng hóa tác động bằng tiền (VNĐ) hoặc phần trăm cải thiện quy đổi được ra tiền; tính tỷ suất hoàn vốn (ROI) dự kiến; mô tả thêm tác động định tính như cải thiện trải nghiệm khách hàng hay giải phóng thời gian nhân viên.
\end{huongdan}

Nội dung của phần này. Chỉ số \tvt{ROI}{Return on Investment -- tỷ suất hoàn vốn} được sử dụng để lượng hóa hiệu quả đầu tư của giải pháp đề xuất.

\subsection{Đóng góp khoa học và hướng mở rộng}
\label{subsec:dong_gop_mo_rong}

\begin{huongdan}
Nêu đóng góp khoa học mới, liên kết trực tiếp với Research Gap; đóng góp thực tiễn cho doanh nghiệp/ngành; hai đến ba hướng nghiên cứu tương lai cụ thể; và tiềm năng công bố (tạp chí/hội nghị phù hợp).
\end{huongdan}

Nội dung của phần này.

\section{Tổng kết chương 3}
\label{sec:tong_ket_c3}

\begin{huongdan}
Điểm lại ngắn gọn: tiêu đề, mục tiêu và câu hỏi nghiên cứu đã xây dựng (Mục~\ref{sec:tieu_de_muc_tieu}, Mục~\ref{sec:rq_finer}); phương pháp và kế hoạch thực hiện (Mục~\ref{sec:phuong_phap_nc}, Mục~\ref{sec:ke_hoach}); kết quả và đóng góp dự kiến (Mục~\ref{sec:ket_qua_du_kien}). Kết thúc bằng câu chuyển ý sang phần Kết luận, ví dụ: toàn bộ đề cương vừa trình bày sẽ được nhìn lại và tự đánh giá ở phần tiếp theo.
\end{huongdan}

Nội dung của phần này.

\end{spacing}
